\documentclass[openany]{ctexbook}
\usepackage{amsmath,amssymb}
\usepackage[hidelinks]{hyperref}
\newcommand{\note}[1]{\footnote{译注：#1}}
\title{代数几何基础第五卷预备笔记\\命题9.3及证明：受约束的超曲面截面}
\author{据用户上传英译编辑稿译出}
\date{2026年11月15日}
\begin{document}
\maketitle
\chapter*{范围与校订状态}
本批连续翻译底本第37页命题9.3的陈述及其跨至第38页的证明，保留原脚注53。下一处为第38页“结合8.18可得推论9.4”的编排提示，不跳过未写出的推论与残缺定理9.5。本批为独立增补，未并入主合订本。底本9.3(c)条件连接词及证明末段的适用范围尚有疑点，译文不将整个命题作为已独立审定的定理；此前8.8缺文等问题均未因此消除。
\chapter*{命题9.3及原稿证明}
\noindent\textbf{命题9.3〔原稿陈述；(c)项措辞待核〕.}
设 $X$ 几何不可约，且 $\dim f(X)\ge2$；取 $x\in X(k)$ 及 $n\ge2$。
\begin{enumerate}
\item[(a)] 考虑 $X$ 相对于 $f_n=u_n f$ 的、经过 $x$ 的超平面截面的线性族。其泛成员所定义的 $Y^{(n)}_\xi$ 几何不可约。
\item[(b)] 设 $L$ 是 $P$ 中经过 $f(x)$ 而不包含 $f(X)$ 的线性子簇。考虑 $X$ 相对于 $f_n$ 的如下超平面截面线性族：它由 $P(n)$ 中“在 $x$ 处与 $L$ 相切”的超平面定义，即由 $P$ 上的 $n$ 次齐次式定义，这些式子限制到 $L$ 后在 $x$ 处至少二阶为零。其泛成员 $Y^{(n)}_\xi$ 几何不可约。\note{这里在 $L$ 上的点应理解为 $f(x)$；沿用底本“在 $x$ 处”的简写。条件是限制到 $L$ 后在该点至少二阶消失，不是沿整个 $L$ 的理想平方条件。}
\item[(c)] 设 $X$ 在 $x$ 处光滑，且 $n\ge3$〔底本接着写：“where $f(X)$ is not contained in a plane defined over $k$”，即“其中 $f(X)$ 不包含于一个在 $k$ 上定义的平面中”；连接关系待核〕。考虑 $X$ 相对于 $f_n$ 的、“在 $x$ 处与 $X$ 相切”的超平面截面 $Y^{(n)}_\xi$ 族。其泛成员所定义的 $Y^{(n)}_\xi$ 几何不可约。\note{底本确实印为“and that $n\ge3$ where…”，不是可确认的“or”。证明随后区分切空间像的维数，并在二维情形使用三次式、较高维情形使用二次式，提示条件可能另有原意；本批不据此猜改连接词，不声称恢复手稿。}
\end{enumerate}

\noindent\textbf{证明（原稿）.}
利用上一节末尾的判据，证明本质上是平凡的。\note{此句为原稿评语，不是本次独立验证的结论。后文照译论证，并单列可核实的记号订正与未决问题。}
在 $P$ 中取包含 $f(x)$ 的仿射坐标开集，并在其逆像上考虑 $X$，\note{底本写“取一个包含 $f(X)$ 的 $P$ 的仿射模型”，一般不能要求整幅射影像包含于一个仿射开集。这里改为经过 $f(x)$ 的仿射坐标邻域，其逆像非空，因 $X$ 不可约而稠密，足以定义下文函数域中的坐标。此局部化并不自行完成对全截面的不可约性论证，后者仍使用上一节判据。}
在(a)中，问题归结为：在坐标 $T_1,T_2,\ldots,T_r$ 中找到三个次数不超过2的多项式 $P,Q,R$，使得
\[
\frac{Q(t)}{P(t)},\qquad \frac{R(t)}{P(t)}
\]
在 $K$ 中于 $k$ 上代数无关。这里 $K$ 是 $X$ 的函数域，$t=(t_1,\ldots,t_r)$ 是由 $T_i$ 定义的 $K$ 中元素系。\note{如 $X$ 未约化，函数域按其不可约约化概形理解。本段坐标可取以 $f(x)$ 为原点，从而后续坐标式在该点消失；这是下文所用的坐标约定。多项式 $P$ 与射影空间 $P$ 同名，沿用底本。}
在(b)中，还要求 $P,Q,R$ 限制到 $L$ 后在 $x$ 处至少二阶消失；此外可设 $L$ 的方程为
\[
T_1=\cdots=T_s=0.
\]
\note{底本这里写“vanish to order two at least on $L$”。若解作 $P,Q,R\in I_L^2$，便与下文 $P=T_1$ 冲突。按(b)项原有括号中的定义，正文明确为限制到 $L$ 后在点上二阶消失；下文 $P,Q,R$ 限制到 $L$ 恒为零，确实满足这一条件。}
最后，(c)中的情形相同，只是 $L$ 为 $X$ 在 $x$ 处切空间的像，并且必要时允许取次数为3的 $P,Q,R$，也就是说，可以稍稍放宽次数。

假设 $\dim f(X)\ge2$ 意味着 $k(t_1,\ldots,t_r)$ 在 $k$ 上的超越次数至少为2。因此，可以找到两个代数无关的坐标函数，不妨记为 $t_1,t_2$。\note{底本印作 $K(t_1,\ldots,t_r)$，但所有 $t_i$ 已在 $K$ 中，故该式等于 $K$，不能表达像的维数；正文订为 $k(t_1,\ldots,t_r)$。}
于是在(a)中取
\[
P=T_1,\qquad Q=T_1^2,\qquad R=T_1T_2.
\]
在(b)中也类似地这样取。注意，此时可选择使 $T_1$ 在 $L$ 上为零的 $t_1$，这是因为 $f(X)\not\subset L$。这意味着存在 $1\le i\le s$ 使 $t_i\ne0$；由于 $t_i$ 在 $x$ 处为零，它不可能是常数，因而不在 $k$ 上代数。\footnote{原脚注53：“〔英译者〕由译者补入。”该脚注在底本中附于“不在 $k$ 上代数”之后。}\note{底本印为 $f(x)\not\subset L^*$，若星号不表示另一个对象，则与(b)中的 $f(x)\in L$ 相抵触；正文按(b)假设改为 $f(X)\not\subset L$，不为无说明的星号添加对象。底本先说存在下标 $i$，随后却写 $t_s$，正文统一为所选 $t_i$。原文紧接着允许设 $k$ 代数闭；在此假设下，函数域中代数于 $k$ 的元素属于 $k$，故所述推理成立。}
（附记：可以假设 $k$ 为代数闭域。）

除去 $f(X)$ 包含在 $X$ 于 $x$ 的切空间经 $f$ 所得的像 $L$ 中这一情形外，(c)由(b)推出。〔若 $\dim X=2$，这一情形确实是例外：一个二次曲面与其相切平面的交迹，一般由两条相交直线组成，因此不可约性不成立。〕\note{方括号为原稿所有。这里保留其例示性质，不扩大为任意特征或任意退化二次曲面的无条件结论。}
但为处理这一情形，在上面明确给出的 $P,Q,R$ 中，可以将 $X$ 换成 $L$ 本身，而在 $L$ 上问题的解是平凡的。\note{这是原稿的归约断言。若仅知 $f(X)\subset L$，在 $L$ 上代数无关的函数限制到 $f(X)$ 后不必仍代数无关。因而本批不把这句视为已补全的论证；需结合原假设与切空间像的维数进一步核查，尤其不能单靠下面的多项式计算确认(c)项的全部范围。}
（若 $\dim L=2$，取
\[
P=T_r^2,\qquad Q=T_r^3,\qquad R=T_r^2T_{r-1},
\]
注意 $Q/P=T_r$ 与 $R/P=T_{r-1}$ 是在 $L$ 上线性无关的线性式，因而代数无关。若 $\dim L\ge3$，则 $T_{r-2},T_{r-1},T_r$ 在 $L$ 上线性无关，取
\[
P=T_r^2,\qquad Q=T_{r-1}^2,\qquad R=T_{r-2}^2.
\]
）\note{底本在这段后缺右括号，正文仅补排版闭合。所取坐标在 $L$ 上的线性无关性按原文保留；在这样的线性空间坐标约定下，两比值的代数无关性可核验，但不自动解决上一条译注所述限制到像的问题。}
\end{document}
